% Options for packages loaded elsewhere
\PassOptionsToPackage{unicode}{hyperref}
\PassOptionsToPackage{hyphens}{url}
%
\documentclass[
]{article}
\usepackage{amsmath,amssymb}
\usepackage{iftex}
\ifPDFTeX
  \usepackage[T1]{fontenc}
  \usepackage[utf8]{inputenc}
  \usepackage{textcomp} % provide euro and other symbols
\else % if luatex or xetex
  \usepackage{unicode-math} % this also loads fontspec
  \defaultfontfeatures{Scale=MatchLowercase}
  \defaultfontfeatures[\rmfamily]{Ligatures=TeX,Scale=1}
\fi
\usepackage{lmodern}
\ifPDFTeX\else
  % xetex/luatex font selection
\fi
% Use upquote if available, for straight quotes in verbatim environments
\IfFileExists{upquote.sty}{\usepackage{upquote}}{}
\IfFileExists{microtype.sty}{% use microtype if available
  \usepackage[]{microtype}
  \UseMicrotypeSet[protrusion]{basicmath} % disable protrusion for tt fonts
}{}
\makeatletter
\@ifundefined{KOMAClassName}{% if non-KOMA class
  \IfFileExists{parskip.sty}{%
    \usepackage{parskip}
  }{% else
    \setlength{\parindent}{0pt}
    \setlength{\parskip}{6pt plus 2pt minus 1pt}}
}{% if KOMA class
  \KOMAoptions{parskip=half}}
\makeatother
\usepackage{xcolor}
\setlength{\emergencystretch}{3em} % prevent overfull lines
\providecommand{\tightlist}{%
  \setlength{\itemsep}{0pt}\setlength{\parskip}{0pt}}
\setcounter{secnumdepth}{-\maxdimen} % remove section numbering
\ifLuaTeX
  \usepackage{selnolig}  % disable illegal ligatures
\fi
\IfFileExists{bookmark.sty}{\usepackage{bookmark}}{\usepackage{hyperref}}
\IfFileExists{xurl.sty}{\usepackage{xurl}}{} % add URL line breaks if available
\urlstyle{same}
\hypersetup{
  hidelinks,
  pdfcreator={LaTeX via pandoc}}

\author{}
\date{}

\begin{document}

\newcommand{\ba}{\begin{aligned}}
\newcommand{\ea}{\end{aligned}}
\newcommand{\bm}{\begin{bmatrix}}
\newcommand{\em}{\end{bmatrix}}
\newcommand{\bc}{\begin{cases}}
\newcommand{\ec}{\end{cases}}
\newcommand{\empty}{\varnothing}
\newcommand{\se}{\subseteq}
\newcommand{\ve}{\varepsilon}
\newcommand{\s}{\sigma}
\newcommand{\0}{\textbf 0}
\newcommand{\1}{\textbf 1}
\newcommand{\I}{\textbf I}
\newcommand{\X}{\textbf X}
\newcommand{\x}{\textbf x}
\newcommand{\y}{\textbf y}
\newcommand{\b}{\textbf b}
\newcommand{\w}{\textbf w}
\newcommand{\P}{\mathcal P}
\newcommand{\S}{\mathcal S}
\newcommand{\T}{\mathcal T}
\newcommand{\A}{\mathcal A}
\newcommand{\B}{\mathcal B}
\newcommand{\C}{\mathcal C}
\newcommand{\L}{\mathcal L}
\newcommand{\c}{\backslash}
\newcommand{\if}{\text{ if }}
\newcommand{\uinf}[1][k]{\bigcup_{{#1} = 1}^{\infty}}
\newcommand{\sinf}[1][k]{\sum_{{#1} = 1}^{\infty}}
\newcommand{\iinf}[1][k]{\bigcap_{{#1} = 1}^{\infty}}
\newcommand{\plim}{\text{plim }}
\newcommand{\var}[2][]{\text{Var}_{#1}\br{#2}}
\newcommand{\e}[2][]{\text{E}_{#1}\br{#2}}
\newcommand{\cov}[3][]{\text{Cov}_{#1}\br{#2,#3}}
\newcommand{\estvar}[1]{\text{Est.Var}\br{#1}}
\newcommand{\estcov}[2]{\text{Est.Cov}\br{#1,#2}}
\newcommand{\asyvar}[1]{\text{Asy.Var}\br{#1}}
\newcommand{\asycov}[2]{\text{Asy.Cov}\br{#1,#2}}
\newcommand{\estasyvar}[1]{\text{Est.Asy.Var}\br{#1}}
\newcommand{\too}{\Longrightarrow}
\newcommand{\sq}{\sqrt}
\newcommand{\dev}[2][i]{({#2}_{#1}-\bar{#2})}
\newcommand{\inv}[1]{{#1}^{-1}}
\newcommand{\add}[1][i]{\sum_{#1 = 1}^n}
\newcommand{\ply}[1][i]{\prod_{#1 = 1}^n}
\newcommand{\br}[1]{\left[#1\right]}
\newcommand{\pr}[1]{\left(#1\right)}
\newcommand{\brace}[1]{\left\{#1\right\}}
\newcommand{\seq}[1][x]{{#1}_1,{#1}_2,\cdots,{#1}_{n}}
\newcommand{\seqadd}[1][x]{{#1}_1+{#1}_2+\cdots+{#1}_{n}}
\newcommand{\dto}{\overset{\text d}{\to}}
\newcommand{\asim}{\overset{\text a}{\sim}}
\newcommand{\st}{\text{ s.t. }}
\newcommand{\om}[1]{\mu^*\pr{#1}}
\newcommand{\lm}[1]{\mu\pr{#1}}
\newcommand{\d}{\text{ d}}
\newcommand{\dm}{\text{ d}\mu}
\newcommand{\dl}{\text{ d}\lambda}
\newcommand{\F}{\mathbb F}
\newcommand{\norm}[1]{\|{#1}\|}
\newcommand{\null}{\text{null }}
\newcommand{\span}{\text{span }}

\textbf{Definition (\(\norm{f}_p\))} Suppose \(\pr{X,\S,\mu}\) is a
measure space, \(0<p<\infty\), and \(f:X\to\F\) is \(\S\)-measurable.
Then the \(p\)-norm of \(f\), denoted \(\norm{f}_p\), is defined by

\[\norm{f}_p = \pr{\int \abs{f}^p\dm}^{1/p}\]

Also when \(p = \infty\), \(\norm{f}_{\infty}\) is called the
\emph{essential supremum} of \(f\), and is defined by

\[\norm{f}_{\infty}:=\inf\brace{t>0:\lm{\brace{x\in X:\abs{f\pr{x}}>t}} = 0}\]

\textbf{Definition (\(L^p\pr{\mu}\))} Suppose \(\pr{X,\S,\mu}\) is a
measure space and \(0<p\le \infty\). The Lebesgue space \(L^p\pr{\mu}\)
is denoted to be the set of all \(\S\)-measurable functions
\(f:X\to \F\) such that

\[\norm{f}_p<\infty\]

\textbf{Remarks} \(L^p\pr{\mu}\) is a vector space. However, it is NOT a
normed vector space. There exists some \(f,g\in L^p\pr{\mu}\) such that
\(f\ne g\) but \(\norm{f-g}_p = 0\).

We construct \emph{equivalence classese}:

\[f\sim g\iff f = g \text{ a.e.}\]

and identify each function with its own equivalence class.

\textbf{Properties of \(L^p\) Spaces}

\textbf{Definition (Dual Exponent)} Suppose \(1\le p\le \infty\). The
\emph{dual exponent} of \(p\) is denoted by \(p'\in \br{1,\infty}\) that
satisfies

\[\frac1p+\frac1{p'}=1\]

\textbf{Examples} \(2' = 2\), \(1'=\infty\), \(\infty' = 1\).

\emph{\textbf{Proposition}}

\begin{itemize}
\item
  (Holder\textquotesingle s Inequality) \(1\le p\le \infty\),
  \(f,h:X\to \F\), then

  \[\norm{fh}_1\le \norm{f}_p\norm{h}_{p'}\]
\item
  \(0<p<q<\infty\), \(\mu\pr{X}<\infty\), then

  \[\norm{f}_p\le \mu\pr{X}^{\pr{q-p}/\pr{pq}}\norm{f}_q\]

  for all \(f\in L^q\pr{\mu}\). Furthermore,
  \(L^q\pr{\mu}\le L^p\pr{\mu}\).
\item
  \(1\le p<\infty\), \(f\in :^p\pr{\mu}\). Then

  \[\norm{f}_p = \sup\brace{\abs{\int fh\dm}:h\in L^{p'}\pr{\mu}\text{ and }\norm{h}_p\le 1}\]
\item
  (Minowsky\textquotesingle s Inequality: \(\norm{\cdot}_p\) is a norm)
  \(1\le p\le \infty\),

  \[\norm{f+g}_p\le \norm{f}_p+\norm{g}_p\]
\end{itemize}

Now that we know \(L^p\) is a normed vector space. We want to show
further that \(L^p\) is a Banach space.

\emph{\textbf{Proposition}} \(1\le p\le \infty\),
\(\brace{f_j}\in L^p\pr{\mu}\) such that for all \(\ve>0\), there exists
\(n\in \N\) such that

\[\norm{f_j-f_k}<\ve,\forall j,k\ge n\]

There exists \(f\in L^p\) such that

\[\lim_{k\to\infty}\norm{f_k-f}_p = 0\]

\emph{\textbf{Proof}} To prove a Cauchy sequence is convergent, it
suffice to show the convergence of any subsequence.

To extract a subsequence with the following properties: Suppose
\(f_0 = 0\),

\[\sinf \norm{f_k-f_{k-1}}<\infty\]

This step is valid because the sequence of Cauchy.

Define \(g_1,g_2,\cdots,g_m\) such that

\[g_m\pr{x} = \sum_{k = 1}^m\abs{f_k\pr{x}-f_{k-1}\pr{x}}\]

and

\[g\pr{x} = \sinf \abs{f_k\pr{x}-f_{k-1}\pr{x}}\]

We have by Triangle Inequality,

\[\norm{g_m}_p\le \sum_{k = 1}^m\norm{f_k\pr{x}-f_{k-1}\pr{x}}_p<\infty\]

Moreover, we have pointwise convergence:

\[g_m\pr{x}\to g\pr{x},\forall x\in X\]

By DCT, we have

\[\ba
&\norm{g}_p = \lim_{m\to\infty}\norm{g_m}_p<\infty\\
\too & g<\infty \text{ a.s}
\ea\]

Therefore,

\[\ba
\sum_{k  =1}^m\pr{f_k\pr{x}-f_{k-1}\pr{x}} & = \lim_{m\to\infty}f_m\pr{x} = f\pr{x}
\ea\]

Moreover, \(\abs{f}\le g\) and \(g\in L^p\), this implies that
\(f\in L^p\).

Finally,
\(\norm{f_k-f}_p\le \liminf_{j\to\infty}\norm{f_k-f_j}_p\le \ve\), for
\(k\) large enough. This implies that

\[\lim_{k\to\infty}\norm{f_k-f}_p=0\]

\end{document}
